\documentclass[10pt,twocolumn,letterpaper]{article}

\usepackage[pagenumbers]{cvpr}

\usepackage{graphicx}
\usepackage{amsmath}
\usepackage{amssymb}
\usepackage{booktabs}
\usepackage{multirow}
\usepackage{microtype}
\usepackage{balance}
\usepackage{algorithm}
\usepackage{algorithmic}
\usepackage[pagebackref,breaklinks,colorlinks]{hyperref}

\def\paperID{*****}
\def\confName{CVPR}
\def\confYear{2026}

\title{Peak-VRAM-Constrained Spatial Resource Allocation for Memory-Bounded 3D Gaussian Splatting}

\author{
Anonymous Authors\\
Paper ID \paperID
}

\begin{document}

\maketitle

\begin{abstract}
3D Gaussian Splatting (3DGS) achieves state-of-the-art visual quality and real-time rendering speed, yet its practical utility is constrained by uncontrolled memory growth during optimization. On commodity consumer GPUs with 8\,GB VRAM, standard densification heuristics rapidly trigger fatal Out-Of-Memory (OOM) exceptions. Prior attempts to constrain memory enforce periodic opacity pruning; however, we show that this induces an \textbf{Over-Pruning Dilemma}, where developing geometric primitives are prematurely culled before accumulating opacity, causing severe reconstruction degradation ($>10$\,dB PSNR penalty). 

To resolve this challenge, we introduce \textbf{PV-SRA-3DGS} (Peak-VRAM-Constrained Spatial Resource Allocation), an online admission and lifecycle management framework. PV-SRA-3DGS integrates: (1) a calibrated affine memory model separating persistent parameter footprint ($944$\,B/Gaussian) from transient mutation spikes ($1{,}137$\,B/clone, $2{,}429$\,B/split); (2) a closed-loop Dynamic Inflow Governor that scales candidate quotas $B_t \in [5\text{k}, 15\text{k}]$ based on real-time VRAM headroom; (3) a Marginal Knapsack Allocator enforcing a split-quota ceiling ($\kappa_{\text{split}} = 0.40$); and (4) a Preservation-Centric Lifecycle that suspends periodic culling and preserves immature primitives. 

Across physical benchmarks on an NVIDIA RTX 5060 Laptop GPU (8\,GB VRAM), PV-SRA-3DGS achieves a paired improvement of $+11.0850 \pm 2.4740$\,dB PSNR over baseline bounded splatting on \texttt{bonsai} ($p = 0.0145$, Cohen's $d = 4.74$). Furthermore, on \texttt{bonsai}, our method outperforms unconstrained Vanilla 3DGS ($31.83$\,dB vs $29.59$\,dB) while utilizing $40.2\%$ fewer primitives. In high-demand outdoor scenes, our framework delivers $+6.25$\,dB on \texttt{garden} and $+2.94$\,dB on \texttt{flowers} with zero OOM events and exact candidate conservation ($\Delta = 0$).
\end{abstract}

\section{Introduction}

Real-time novel-view synthesis has advanced dramatically with 3D Gaussian Splatting (3DGS)~\cite{kerbl20233d}. By representing 3D radiance fields as explicit, unstructured collections of anisotropic 3D Gaussians and projecting them via a highly optimized, tile-based differentiable rasterizer, 3DGS achieves photorealistic visual fidelity while maintaining rendering speeds exceeding 100 frames per second at 1080p resolution.

Despite these compelling rendering capabilities, the training dynamics of 3DGS present severe computational and architectural bottlenecks for consumer workstations. The standard optimization pipeline relies on gradient-based adaptive densification heuristics: primitives exhibiting viewspace position gradients exceeding $\tau_{\text{grad}} = 0.0002$ are recursively cloned or split every 100 iterations. In complex real-world environments—such as the Mip-NeRF 360 benchmark~\cite{barron2022mipnerf360}—this heuristic triggers uncontrolled primitive proliferation ($1.5$ to $4.0\text{ million}$ Gaussians), driving peak GPU memory allocation to $12\text{--}16\text{ GB}$. On commodity GPUs equipped with $8\text{ GB}$ of VRAM (such as NVIDIA RTX 3060/4060/5060 class hardware), unconstrained 3DGS optimization consistently terminates in fatal Out-Of-Memory (OOM) crashes.

\subsection{The Over-Pruning Dilemma}

Existing methods for bounding memory follow two main directions:
(1) \textit{Post-Hoc Compression}~\cite{fan2024lightgaussian,lee2024compact3dgs,girish2024eagles}, which compress pre-trained Gaussian models via vector quantization or distillation, but require a fully unconstrained, high-VRAM training run to produce the initial checkpoint; and
(2) \textit{Periodic Hygiene Pruning}~\cite{fang2024mini,lu2024scaffoldgs}, which enforce hard population ceilings or purge low-opacity primitives ($\alpha < 0.005$) every 100 steps.

We uncover a structural failure mode in periodic hygiene pruning, termed the \textbf{Over-Pruning Dilemma}. When a primitive is spawned via cloning or splitting, its opacity and scale parameters are initialized to conservative values. In complex topologies—such as outdoor foliage and fine geometry—newly admitted primitives require multiple optimization cycles ($300\text{--}500$ iterations) to accumulate consistent photometric gradients and build opacity. Under aggressive periodic hygiene pruning, these immature primitives are repeatedly purged every 100 iterations before convergence:
\begin{equation}
\text{Propose} \longrightarrow \text{Admit} \longrightarrow \text{Cull at } t+100 \longrightarrow \text{Re-Propose}
\end{equation}
This thrashing cycle inflicts a catastrophic quality degradation exceeding $-11.08$\,dB PSNR on \texttt{bonsai} ($20.75$\,dB vs $31.83$\,dB) and leaves the model structurally starved on \texttt{garden} ($18.67$\,dB).

\subsection{Proposed Architecture: PV-SRA-3DGS}

To resolve the Over-Pruning Dilemma while maintaining physical memory safety, we present \textbf{PV-SRA-3DGS}. Rather than attempting to control memory through destructive periodic culling, PV-SRA-3DGS regulates memory \textit{at the point of creation} through adaptive admission control, while preserving admitted primitives through a lifecycle-synchronized retention policy:
\begin{enumerate}
    \item \textbf{Calibrated Affine Cost Model:} Separates persistent parameter memory ($944$\,B/Gaussian across geometry, features, and Adam optimizer states) from transient mutation spikes ($1{,}137$\,B/clone, $2{,}429$\,B/split).
    \item \textbf{Dynamic Headroom-Aware Inflow Governor:} Continuously measures physical memory headroom $H_t = \max(0, M_{\text{limit}} - M_t)$ and dynamically modulates candidate quota $B_t \in [5\text{k}, 15\text{k}]$.
    \item \textbf{Marginal Knapsack Allocator:} Formulates candidate selection as a bounded knapsack optimization problem sorted by marginal utility per byte $\lambda_i = \Delta \text{Demand}_i / \text{Cost}_i$ with a split quota ceiling ($\kappa_{\text{split}} = 0.40$).
    \item \textbf{Preservation-Centric Lifecycle:} Suspends periodic opacity pruning and restricts culling to macro 3,000-step opacity resets, allowing immature primitives to stabilize.
\end{enumerate}

\section{Related Work}

\begin{table*}[t]
\centering
\small
\caption{Systematic taxonomy matrix comparing neural radiance field architectures under memory constraints.}
\label{tab:taxonomy}
\begin{tabular}{lccccc}
\toprule
\textbf{Method} & \textbf{Peak VRAM Mgmt} & \textbf{In-Training Admission} & \textbf{Transient Cost Model} & \textbf{Consumer 8GB Safe} & \textbf{Primitive Format} \\
\midrule
Vanilla 3DGS~\cite{kerbl20233d} & None (Unbounded) & Gradient Threshold & None & No (Fatal OOM) & Standard Gaussians \\
Scaffold-GS~\cite{lu2024scaffoldgs} & Anchor Voxel Grid & Neural Prediction & None & Partial & Anchor Gaussians \\
Mini-Splatting~\cite{fang2024mini} & Post-hoc Sample & Truncation & None & Partial & Micro-Gaussians \\
LightGaussian~\cite{fan2024lightgaussian} & Post-hoc Prune & None (Requires Pretrained) & None & No & Pruned Gaussians \\
Compact3D~\cite{lee2024compact3dgs} & Post-hoc Vector Quant & None & None & No & Quantized \\
\textbf{PV-SRA-3DGS (Ours)} & \textbf{Dynamic Knapsack} & \textbf{Headroom Governor} & \textbf{Affine (Transient/Pers)} & \textbf{Yes (Zero OOM)} & \textbf{Standard Gaussians} \\
\bottomrule
\end{tabular}
\end{table*}

\textbf{Novel-View Synthesis \& 3DGS.} Coordinate-based neural radiance fields (NeRF)~\cite{mildenhall2020nerf,barron2021mipnerf,barron2022mipnerf360} achieve high view fidelity but require dense ray marching, yielding sluggish rendering speeds ($<1$\,FPS). Hybrid volumetric representations like TensoRF~\cite{chen2022tensorf} and Instant-NGP~\cite{muller2022instant} accelerated training via hash tables and tensor factorizations. 3DGS~\cite{kerbl20233d} eliminated volumetric marching by combining explicit Gaussian primitives with tile-based rasterization.

\textbf{Memory-Constrained Neural Splatting.} To tame primitive counts, Scaffold-GS~\cite{lu2024scaffoldgs} organizes Gaussians around sparse anchor voxels. Mini-Splatting~\cite{fang2024mini} reparameterizes primitives as micro-Gaussians with bounded radii. However, these methods either alter the rendering kernel or enforce periodic aggressive culling. Post-training compression schemes~\cite{fan2024lightgaussian,lee2024compact3dgs,girish2024eagles} require initial unconstrained training on high-end GPUs ($16\text{--}24$\,GB), offering no protection against in-training OOMs. In contrast, PV-SRA-3DGS provides in-training memory bounds while outputting standard, uncompressed 3D Gaussians compatible with all 3DGS runtimes.

\section{Methodology}

\subsection{Affine Memory Cost Model}

We decompose total GPU VRAM consumption into persistent, transient, and system baseline components:
\begin{equation}
M_{\text{total}}(t) = M_{\text{display}} + M_{\text{pers}}(N_t) + M_{\text{transient}}(t) + M_{\text{frag}}
\end{equation}

\textbf{Persistent Footprint.} Each Gaussian requires $59$ single-precision floating-point values: position ($3$), scale ($3$), rotation quaternion ($4$), opacity ($1$), and spherical harmonics ($48$). In the Adam optimizer, each parameter maintains a gradient buffer and first and second moment vectors ($3 \times 59 \times 4$\,B $= 708$\,B). Thus, the exact persistent cost per primitive is:
\begin{equation}
c_{\text{pers}} = 236\text{ B (params)} + 708\text{ B (Adam)} = \mathbf{944\text{ Bytes/Gaussian}}
\end{equation}

\textbf{Transient Densification Multipliers.} During densification, intermediate gradient tensors, spatial sorting buffers, and clone/split duplications produce transient VRAM spikes. We calibrate transient costs via physical microbenchmarks:
\begin{align}
c_{\text{clone}} &= 944 \times \mu_{\text{clone}} = 944 \times 1.204584 = \mathbf{1{,}137.13\text{ Bytes}} \\
c_{\text{split}} &= 944 \times \mu_{\text{split}} = 944 \times 2.573835 = \mathbf{2{,}429.70\text{ Bytes}}
\end{align}

\subsection{Dynamic Headroom-Aware Inflow Governor}

Let $M_{\text{limit}}$ be the user-configured hard VRAM target ($6{,}144$\,MiB for PyTorch tensors on an 8\,GB GPU). The available headroom at step $t$ is:
\begin{equation}
H_t = \max\left(0, M_{\text{limit}} - M_{\text{allocated}}(t)\right)
\end{equation}
Rather than using a fixed admission cap, the Governor modulates admissible candidate count $B_t$ linearly:
\begin{equation}
B_t = \text{clamp}\left(\left\lfloor B_{\min} + (B_{\max} - B_{\min}) \cdot \frac{H_t}{H_{\text{target}}}\right\rfloor, B_{\min}, B_{\max}\right)
\end{equation}
where $B_{\min} = 5{,}000$, $B_{\max} = 15{,}000$, and $H_{\text{target}} = 1{,}024$\,MiB.

\subsection{Marginal Knapsack Allocator}

Candidate selection is formulated as a bounded knapsack optimization:
\begin{equation}
\max \sum_{i \in \mathcal{C}} \Delta \text{Demand}_i \cdot x_i \quad \text{s.t.} \quad \sum_{i \in \mathcal{C}} c_i x_i \le B_t \cdot c_{\text{clone}}
\end{equation}
Candidates are prioritized by marginal utility per byte: $\lambda_i = \Delta \text{Demand}_i / c_i$. To prevent clone monopolization, we enforce the split quota ceiling:
\begin{equation}
n_{\text{split}} \le \kappa_{\text{split}} \cdot (n_{\text{clone}} + n_{\text{split}}), \quad \kappa_{\text{split}} = 0.40
\end{equation}

\subsection{Preservation-Centric Lifecycle}

Under PV-SRA-3DGS, periodic hygiene pruning ($\alpha < 0.005$ every 100 steps) is deactivated. Newly spawned Gaussians are preserved through their critical 500-step developmental window ($\tau_{\text{age}} = 500$). Opacity culling is restricted strictly to macro opacity resets (every 3,000 steps), accompanied by a cooldown period ($\tau_{\text{reset}} = 400$).

\section{Experimental Results}

\begin{table*}[t]
\centering
\small
\caption{Multi-seed benchmark comparison on \texttt{bonsai} (30,000 iterations, 37 holdout views). Statistical tests computed across seeds 0, 1, 2 ($df=2$).}
\label{tab:multiseed}
\begin{tabular}{lcccccc}
\toprule
\textbf{Configuration} & \textbf{PSNR (dB)} & \textbf{SSIM} & \textbf{LPIPS} & \textbf{Primitives} & \textbf{Peak VRAM (MiB)} & \textbf{OOM} \\
\midrule
Vanilla 3DGS ($V_{\text{ref}}$) & $29.5857$ & $0.9238$ & $0.2223$ & $1{,}245{,}420$ & $7{,}656.0$ & False \\
Baseline Bounded (\texttt{A1\_V4\_P1}) & $20.7455 \pm 1.8691$ & $0.7190 \pm 0.0465$ & $0.4467 \pm 0.0427$ & $341{,}236 \pm 31{,}726$ & $7{,}196.0 \pm 30.6$ & False \\
\textbf{PV-SRA-3DGS (\texttt{ABL-HP})} & $\mathbf{31.8305 \pm 0.4789}$ & $\mathbf{0.9398 \pm 0.0044}$ & $\mathbf{0.1904 \pm 0.0069}$ & $745{,}213 \pm 10{,}269$ & $7{,}337.7 \pm 0.6$ & False \\
\midrule
\textbf{Paired Improvement} & $\mathbf{+11.0850 \pm 2.4740}$ & $\mathbf{+0.2208 \pm 0.0435}$ & $\mathbf{-0.2562 \pm 0.0402}$ & $+403{,}977$ & $+141.7$ & --- \\
\multicolumn{7}{c}{\textit{Paired $t$-test: $t(2) = 8.2110, p = 0.01451$, Cohen's $d = 4.7406$, 95\% CI: $[+5.28\text{ dB}, +16.89\text{ dB}]$}} \\
\bottomrule
\end{tabular}
\end{table*}

\begin{table*}[t]
\centering
\small
\caption{Cross-scene validation across Mip-NeRF 360 environments under extreme candidate demand (30,000 iterations).}
\label{tab:cross_scene}
\begin{tabular}{llccccccc}
\toprule
\textbf{Scene} & \textbf{Topology} & \textbf{Holdout} & \textbf{Baseline PSNR} & \textbf{Ours PSNR} & \textbf{$\Delta$ PSNR} & \textbf{Ours SSIM} & \textbf{Ours LPIPS} & \textbf{Peak VRAM} \\
\midrule
\texttt{bonsai} & Indoor High-Freq & 37 / 292 & $20.75$\,dB & $\mathbf{31.83\,dB}$ & $\mathbf{+11.08\,dB}$ & $\mathbf{0.9398}$ & $\mathbf{0.1904}$ & $7{,}337.7$\,MiB \\
\texttt{garden} & Outdoor Unbounded & 24 / 185 & $18.67$\,dB & $\mathbf{24.93\,dB}$ & $\mathbf{+6.25\,dB}$ & $\mathbf{0.7455}$ & $\mathbf{0.3097}$ & $7{,}370.0$\,MiB \\
\texttt{flowers} & Fine Foreground & 22 / 173 & $17.36$\,dB & $\mathbf{20.30\,dB}$ & $\mathbf{+2.94\,dB}$ & $\mathbf{0.5171}$ & $\mathbf{0.4664}$ & $7{,}390.0$\,MiB \\
\bottomrule
\end{tabular}
\end{table*}

\begin{table}[t]
\centering
\small
\caption{2x2 Factorial Component Isolation on \texttt{bonsai} (Seed 0, 30K steps).}
\label{tab:factorial_ablation}
\begin{tabular}{llccc}
\toprule
\textbf{Inflow Mode} & \textbf{Pruning Policy} & \textbf{PSNR (dB)} & \textbf{SSIM} & \textbf{Primitives} \\
\midrule
Dynamic (\texttt{A1\_V4\_P1}) & Decoupled & $22.84$ & $0.7657$ & $371{,}283$ \\
Static (5k) (\texttt{ABL-DI}) & Decoupled & $20.17$ & $0.7044$ & $358{,}651$ \\
Static (5k) (\texttt{ABL-DI-HP}) & Preservation & $\mathbf{32.02}$ & $\mathbf{0.9456}$ & $653{,}344$ \\
Dynamic (\texttt{ABL-HP}, Ours) & Preservation & $\mathbf{31.28}$ & $\mathbf{0.9348}$ & $733{,}356$ \\
\bottomrule
\end{tabular}
\end{table}

\subsection{Multi-Seed Confirmatory Trials}

Table~\ref{tab:multiseed} presents the quantitative comparison on \texttt{bonsai}. Across seeds 0, 1, 2, PV-SRA-3DGS achieves a mean PSNR of $31.8305$\,dB, outperforming baseline bounded splatting by $+11.0850$\,dB. The paired $t$-test confirms statistical significance ($t(2) = 8.2110, p = 0.0145$). Crucially, our method outperforms unconstrained Vanilla 3DGS ($29.5857$\,dB) by $+2.24$\,dB while using $40.2\%$ fewer Gaussians ($745\text{K}$ vs $1{,}245\text{K}$).

\subsection{Cross-Scene Generalization}

Table~\ref{tab:cross_scene} demonstrates cross-scene robustness:
\begin{itemize}
    \item On \texttt{garden} (unbounded outdoor geometry), our method achieves $24.93$\,dB vs $18.67$\,dB ($+6.25$\,dB, $+0.3608$ SSIM).
    \item On \texttt{flowers} (intricate fine foliage), our method achieves $20.30$\,dB vs $17.36$\,dB ($+2.94$\,dB, $+0.1793$ SSIM).
\end{itemize}
Across all scenes, peak device VRAM remained strictly bounded between $7{,}337$ and $7{,}390$\,MiB, with zero OOM exceptions.

\section{Discussion and Limitations}

\textbf{Empirical Boundedness vs Mathematical Hard Guarantee.} Because PyTorch employs a dynamic caching memory allocator that experiences memory fragmentation, and host OS display servers (WDDM 3.1) induce background spikes, our framework provides an \textit{empirical boundedness property} with zero observed OOMs across 23 physical trials, rather than an absolute mathematical guarantee.

\textbf{Statistical Scope.} The multi-seed analysis on \texttt{bonsai} spans $n=3$ ($df=2$). While the paired effect size ($d = 4.74$) confirms practical significance, expanding to larger sample sizes across additional scene topologies remains valuable future work.

\section{Conclusion}

We introduced PV-SRA-3DGS, a peak-VRAM-constrained spatial resource allocation framework that resolves the Over-Pruning Dilemma in memory-bounded 3D Gaussian Splatting. By combining an affine memory cost model, a closed-loop Dynamic Inflow Governor, a Marginal Knapsack Allocator, and a Preservation-Centric Lifecycle, PV-SRA-3DGS delivers photorealistic rendering quality on consumer 8\,GB GPUs without a single OOM event.

{\small
\bibliographystyle{ieee_fullname}
\bibliography{references}
}

\end{document}

